% =============================================================================
% SOBRE LOS AUTORES
% =============================================================================

\chapter*{Sobre los Autores}
\addcontentsline{toc}{chapter}{Sobre los Autores}

\bigskip

% -----------------------------------------------------------------------------
\subsection*{Froylan Jiménez Sánchez}
% -----------------------------------------------------------------------------

\noindent
\textbf{Froylan Jiménez Sánchez} es docente de matemáticas y física con más de diez años
de experiencia en educación superior. Es Especialista en Matemática Avanzada y Magíster en
Estadística Aplicada, con investigación centrada en la intersección entre métodos
estadísticos, inteligencia artificial y sistemas agroindustriales.

\medskip

Se desempeña como docente e investigador vinculado al \textbf{Grupo PROPED}
(\textit{Grupo de Investigación en Procesos Pedagógicos y Didácticos}) de la
\textit{Universidad de Sucre} (Sincelejo, Colombia), institución donde integra el
análisis estadístico y computacional a la formación cuantitativa de estudiantes de
pregrado en ciencias naturales, ingenierías y ciencias de la salud.

\medskip

Sus áreas de investigación comprenden el aprendizaje automático (\textit{Machine Learning})
aplicado a la agricultura de precisión, la bioimpedancia espectroscópica para la evaluación
no destructiva de madurez en aguacate Hass, los \textit{Fuzzy Cognitive Maps} como
herramienta de modelado de sistemas complejos, y la estadística aplicada al análisis de
datos agronómicos y educativos.

\medskip

Entre sus contribuciones científicas más recientes se destacan un artículo publicado en
\textbf{MDPI \textit{Computers}} (2026) sobre la evaluación de la madurez del aguacate
Hass mediante bioimpedancia espectroscópica y mapas cognitivos difusos, así como un
capítulo de libro en editorial \textbf{Springer} sobre aplicaciones de Machine Learning
en sistemas agrícolas.

\bigskip\bigskip

% -----------------------------------------------------------------------------
\subsection*{Liliana Margarita Vitola Garrido}
% -----------------------------------------------------------------------------

\noindent
\textbf{Liliana Margarita Vitola Garrido} es Licenciada en Matemáticas de la
\textit{Universidad de Sucre} (2004) y Magíster en Estadística Aplicada de la
\textit{Fundación Universidad del Norte} (Barranquilla, 2012). Su trayectoria
académica combina una formación matemática rigurosa con una especialización en
métodos estadísticos orientados a la investigación aplicada.

\medskip

Actualmente se desempeña como \textbf{Docente de Estadística, Diseño de Experimentos
y Matemáticas} en la \textit{Universidad de Sucre} (Sincelejo, Colombia), institución
en la que ha contribuido a la formación cuantitativa de múltiples generaciones de
profesionales de las ciencias naturales, la salud y las ingenierías. Su labor docente
se caracteriza por vincular los fundamentos teóricos de la estadística con su
aplicación directa a problemáticas regionales concretas.

\medskip

Sus áreas de investigación incluyen la regresión logística aplicada al análisis del
rendimiento académico, la lógica difusa (\textit{Fuzzy Logic}) y las funciones de
probabilidad, la validación de instrumentos de medición en salud mental, y el
tratamiento de agua mediante coagulantes naturales como alternativa para comunidades
con acceso limitado a saneamiento básico.

\medskip

Su producción científica abarca publicaciones en revistas indexadas nacionales e
internacionales. Entre las más representativas figuran: \textit{``Fuzzy Logic and
Probability Functions''}, publicado en el \textbf{Colombia International Journal of
Mechanical Engineering and Technology} (2019); \textit{``Supply and Demand for
Drinking Water in Sincelejo City, Colombia: A Review of Alternative Solutions''},
en \textbf{IAEME Publication} (2019); \textit{``Regresión Logística: Una aplicación
en la identificación de variables que inciden en el rendimiento académico en el área
de matemáticas''}, editado por la \textbf{Universidad Militar Nueva Granada} (2015);
y \textit{``Logistic Regression: For the Identification of Socio-Economic Variables
that Influence on the Academic Performance of Students''}, en el
\textbf{Indian Journal of Science and Technology} (2017). Completa su perfil
investigativo un trabajo sobre la actividad coagulante de semillas de
\textit{Moringa oleifera} para el tratamiento de aguas residuales, que evidencia
su interés en soluciones ambientales de bajo costo para el Caribe colombiano.
